% Generado por tablas_plan3.py a partir de plan3/salidas/. No editar a mano.
% Segundo estudio, llegadas aleatorias: casos en que el modelo citado recibe cada etiqueta citable frente
% a cada una de las otras referencias (la misma regla, con esa referencia en el lugar del fijo ajustado al
% segundo).
% fuente: plan3/salidas/veredicto3.md, tabla 'Referencias al lado' (filas P1-P6 con llegadas aleatorias).
% a: mejora en una sola tanda. Entorno: \setstretch{1}, \footnotesize y \setlength{\tabcolsep}{3pt}.
\begin{tabular}{>{\raggedright\arraybackslash}p{4.2cm}>{\raggedright\arraybackslash}p{3.0cm}>{\raggedright\arraybackslash}p{1.8cm}>{\raggedright\arraybackslash}p{3.0cm}>{\raggedright\arraybackslash}p{1.8cm}}
\toprule
Frente a & FUNCIONA & MEJORA & EMPATA & FALLA \\
\midrule
Actuado ajustado & P2, P6 & P1 & P4, P5 & P3 \\
Actuado por defecto & P1 a P6 & --- & --- & --- \\
Fijo ajustado desde 10~s & P1, P2, P6 & P4 & P3\textsuperscript{a}, P5 & --- \\
Fijo de verdes de 5 en 5 & P1, P2, P6 & --- & P4, P5 & P3 \\
Fijo por franja & P1, P2, P6 & --- & P4, P5 & P3 \\
Regla simple & P1, P2, P6 & --- & P3, P4\textsuperscript{a}, P5 & --- \\
Regla reactiva & P1 a P6 & --- & --- & --- \\
Max-Pressure & P1 a P6 & --- & --- & --- \\
Ciclo mínimo & P1, P2, P6 & --- & P3, P4\textsuperscript{a}, P5 & --- \\
\bottomrule
\end{tabular}
